\chapter{The Main Structural Theorem, GRH as Corollary~1, and the Classical Automorphic--Galois Interface}
\label{ch:01}

In this chapter we place the entire proof on its proper structural foundation. Rather than presenting the Generalized Riemann Hypothesis as the end of an opaque chain of ad~hoc estimates, we isolate a single geometric-arithmetic theorem—\cref{thm:01:main-structural}, the \emph{Main Structural Theorem for Arakelov-Filtered Automorphic Period Towers}—and deduce the Generalized Riemann Hypothesis for all Dirichlet $L$-functions (and the Riemann Hypothesis for $\zeta(s)$) immediately as \cref{cor:01:grh}. We then carry out the elementary reductions to primitive characters and the right half-plane $\Re(s) > 1/2$ (\cref{prop:01:imprimitive}), and finally formulate and verify the \emph{Classical Automorphic--Galois Interface} (\cref{thm:01:classical-automorphic-galois}), replacing the artificially weakened Galois-representation and local-parameter surrogates of the formal archive with exact theorems from the published literature.

\section{Arakelov-Filtered Automorphic Period Towers and the Main Structural Theorem}
\label{sec:01:main-structural}

Fix a primitive Dirichlet character $\psi$ of modulus $M \ge 1$ and order $o = \ord(\psi) \ge 1$ (with $M = 1$, $o = 1$, $\psi = \mathbf{1}$ corresponding to the Riemann zeta function $\zeta(s)$). Let $\Xi = \langle \psi \rangle$ be the cyclic group generated by $\psi$, and define the \emph{positively weighted cyclic ramp sum}
\begin{equation}
\label{eq:01:gpsi-def}
  G_\psi(y) \;:=\; G(\mathbf{1}, y) \;+\; \frac{1}{3o + 1} \sum_{\substack{\xi \in \Xi \\ \xi \neq \mathbf{1}}} \Re G(\xi, y) \qquad (y \ge 1),
\end{equation}
where for any admissible Dirichlet character $\chi \pmod{N}$ (meaning $N = 1$, or $N \ge 2$ and $\chi \neq \mathbf{1}$), the centered prime-power ramp sum is
\begin{equation}
\label{eq:01:gchi-def}
  G(\chi, y) \;:=\;
  \begin{cases}
    \displaystyle \sum_{1 \le n \le y} \Lambda(n)\!\left(\frac{y}{n} - 1\right) - y\log y + (1 + \gamma)y, & \chi = \mathbf{1}\ (N = 1), \\[10pt]
    \displaystyle \sum_{1 \le n \le y} \Lambda(n)\chi(n)\!\left(\frac{y}{n} - 1\right) + y\,\frac{L'}{L}(1, \chi), & \chi \neq \mathbf{1}.
  \end{cases}
\end{equation}
Here $\Lambda$ is the von~Mangoldt function and $\gamma$ is Euler's constant. The weight $(3o+1)^{-1}$ in \eqref{eq:01:gpsi-def} is completely natural: it arises directly from the biased-Bernoulli coset distribution on the Galois group of an auxiliary CM field $E/\QQ$ (\cref{ch:03,ch:04}), where the trivial coset receives weight $4/(3o+1)$ and each of the $o-1$ nontrivial cosets receives weight $3/(3o+1)$, so subtracting the balanced base leaves excess weight $1$ on the trivial character and $(3o+1)^{-1}$ on every $\xi \in \Xi \setminus \{\mathbf{1}\}$. Crucially, all weights in \eqref{eq:01:gpsi-def} are strictly positive, ensuring that the poles of the symmetrized Laplace transform of $G_\psi$ at zeros of $L(s,\xi)$ never cancel (\cref{ch:02}).

To state the Main Structural Theorem cleanly, we package the geometric engine of Parts~II--VI into a single structural concept.

\begin{definition}[Arakelov-Filtered Automorphic Period Tower]
\label{def:01:period-tower}
Let $0 < \varepsilon < 1/52$, let $x \ge 2$ be a scale parameter, and let $E/\QQ$ be a finite abelian CM field with $\Gal(E/\QQ) = G = (\ZZ/2\ZZ)^2 \times A$ containing an imaginary quadratic subfield $K = \QQ(\sqrt{-d})$ (with $\Phi = \Gal(E/K)$) and realizing the character group $\Xi = \langle \psi \rangle$ on $A / A_0$ via a twisted Artin surjection (\cref{sec:03:canonical-frame}). Set $n_0 = |A_0| = 2^e$, $L = x^\varepsilon$, $b = \exp(x^\varepsilon)$, and $a = x b$. An \emph{Arakelov-filtered automorphic period tower} of length $m$ and base rank $k$ over $E$ (with odd step increment $t$, even layer ranks $N_j = k + j(t+1)$ for $0 \le j \le m$, a Behrend-thick subset $T_2 \subseteq A_0$, and a balanced top subset $\TT \subseteq A$) consists of:
\begin{enumerate}[label=\textup{(\roman*)}]
\item a growing sequence of positive-definite Hermitian $\OO_E$-lattices $(V_j, \Lambda_j)$ of ranks $N_j = k + j(t+1)$ ($0 \le j \le m$), equipped with definite unitary theta-lift kernels and iterated seesaw period evaluations into a dual tower of automorphic row lattices $\mathcal{R}_u$ ($u \in \TT$) of rank $D_{\mathrm{all}} > 0$ (\cref{ch:07,ch:11});
\item an \emph{inert-prime height-$2$ filtration} at every prime $p \le a$ inert in $K/\QQ$, generated by Lubin--Tate formal group derivatives and cyclic orbital linear-programming bounds (\cref{ch:07,ch:08}), glued across all finite places into an Arakelov-normalized Hermitian lattice $\mathcal{M}_{\mathrm{gl}}$ (\cref{ch:09}) whose total non-inert deficit across all four remaining prime classes is bounded with exponent $e_L = 0$ (\cref{ch:10});
\item an \emph{upper covolume bound} on the dual row lattice $\mathcal{R}_u^\vee$ deduced via seesaw periods, analytic kernel estimates, signature-$(p,q)$ Shimura flag-bundle mean heights, and deep covering descent (\cref{ch:11,ch:12,ch:13,ch:14});
\item a \emph{lower spectral frame bound} (near-orthogonality and quantitative limit multiplicity) deduced from four-way projected Poincar\'e unfolding, rotated wave-grid diagonal bounds, Gelfand--Graev off-diagonal decay, and definite-side closure at all primes including $\ell \in \{2,3\}$ (\cref{ch:15,ch:16,ch:17}); and
\item a \emph{Galois-cohomological integrality bound} controlling twisted-place integrality losses via the higher Tate--Shafarevich length bound $\le \sqrt{b} = \exp(\frac{1}{2}x^\varepsilon)$ across all odd-weight layers up to cutoff $t \le x^{400}$ (\cref{ch:18,ch:19,ch:20}).
\end{enumerate}
\end{definition}

The entire architectural achievement of this monograph is encapsulated in the following theorem.

\begin{theorem}[Main Structural Theorem: The Stepanov--Bost Slope Principle]
\label{thm:01:main-structural}
Let $\psi$ be any primitive Dirichlet character of modulus $M \ge 1$ and order $o \ge 1$, and fix any $\varepsilon \in (0, 1/52)$. Then there exists an explicit threshold $Y_0 = Y_0(\psi, \varepsilon) \ge 2$ such that for every $Y_{\min} \ge Y_0$, there exists a scale $Y \ge Y_{\min}$ for which the following three structural assertions hold simultaneously:
\begin{enumerate}[label=\textup{(\alph*)}]
\item \emph{(Existence of the Arakelov-Filtered Period Tower.)} For every $x \in [Y, Y^{3/2}]$, there exist an auxiliary prime pair $(q, d)$ with $q \equiv 3 \pmod 8$, $d \equiv 7 \pmod{16}$, $(\frac{-d}{q}) = 1$, a host abelian CM field $E/\QQ$ realizing $\langle \psi \rangle$, a Behrend-thick set $T_2 \subseteq A_0$ of density $\delta_2 = |T_2|/n_0 \ge x^{-\varepsilon}$, a balanced top set $\TT \subseteq A$ of density $\delta_1 \asymp x^{-\varepsilon/8}$, and an Arakelov-filtered automorphic period tower (\cref{def:01:period-tower}) at the minimax tower parameters
\begin{equation}
\label{eq:01:minimax-params}
  n_0 \in \bigl[x^{2+8\varepsilon},\, 2x^{2+8\varepsilon}\bigr], \qquad
  m \in \bigl[x^{1/2-\varepsilon/16},\, 2x^{1/2-\varepsilon/16}\bigr], \qquad
  k \in \bigl[x^{9/2+10\varepsilon},\, 2x^{9/2+10\varepsilon}\bigr],
\end{equation}
with $m, k$ even, $t \le x^{400}$ odd, and nonvanishing period matrices ($D_{\mathrm{all}} > 0$, $\rk(R) > 0$).

\item \emph{(Pointwise Main Inequality $(\star)$ via the Two-Sided Slope Sandwich.)} At every $x \in [Y, Y^{3/2}]$ and every translate $u_0 \in A_0$, applying the sharpened multi-map Bost slope inequality (\cref{thm:05:bost-slope}) and the Paley--Zygmund frame codimension inequality (\cref{thm:05:frame-codim}) simultaneously to the glued lattice $\mathcal{M}_{\mathrm{gl}}$ and the dual row lattice $\mathcal{R}^\vee$ cancels the unknown archimedean period slope $\mu(R R^\vee)$ identically and yields the Pointwise Main Inequality $(\star)$ with locked constant $\mathrm{CErr} = e_D + e_L + 3 = 5$:
\begin{equation}
\label{eq:01:star-preview}
  \frac{1}{|\TT|}\sum_{u \in \TT} \bigl(\ASK(x, u, u_0) - \MAIN(x, u)\bigr)
  \;\le\; x^{5\varepsilon}\,\mathcal{E}_{\mathrm{br}}(x, n_0, m, k),
\end{equation}
where the bracketed error $\mathcal{E}_{\mathrm{br}}$ is the sum of the four competing structural scales $x/m$, $xm/k$, $m$, and $k/n_0^2$ multiplied by $n_0 m b x^{O(\varepsilon)}$.

\item \emph{(Emergence of the Square-Root Ramp Bound.)} Averaging \eqref{eq:01:star-preview} over $u_0 \in A_0$, expanding $\ASK - \MAIN$ via the Balance Identity (\cref{thm:03:balance-identity}), evaluating the four competing scales at the minimax tower length $m \asymp x^{1/2-\varepsilon/16}$ in \eqref{eq:01:minimax-params}, and dividing by the thick-set density $\delta_2 \ge x^{-\varepsilon}$ yields the uniform one-sided ramp-sum bound
\begin{equation}
\label{eq:01:gpsi-bound}
  G_\psi(x) \;\le\; x^{1/2 + 9\varepsilon} \qquad \text{for all } x \in \bigl[Y,\, Y^{3/2}\bigr].
\end{equation}
\end{enumerate}
\end{theorem}
\leanref{src/Development/AutomorphicArithmetic/CharacterBoxes/ZeroFreeAllOfTopG.lean:26}{UpperTheory.Top.zeroFreeAll\_of\_topG}
\leanref{src/Development/AutomorphicArithmetic/ThetaCorrespondence/ThetaPairLifts/ThirdPlaceSeries/TopPlug.lean:12}{UpperTheory.Top.zero\_free\_all\_of\_green\_v\_win}

\begin{proof}[Architecture of the proof of \cref{thm:01:main-structural}]
Parts~(a), (b), and (c) assemble the six parts of this monograph in strict acyclic order:
\begin{enumerate}[label=\textup{\arabic*.}]
\item The \emph{Classical Automorphic--Galois Interface} (\cref{thm:01:classical-automorphic-galois} below) feeds into Part~VI (\cref{ch:18,ch:19,ch:20}), where the integrated Selmer filtration, Kolyvagin/Stickelberger minus-part bound, four-step cohomological bound for $\GL_N(\OO_L)$, $2$-power Bernoulli units, and graded deformation congruences establish the unconditional Tate--Shafarevich window bound $\le \exp(\frac{1}{2}x^\varepsilon)$ on arbitrarily large exponential windows $[W, W^2]$ at cutoff $401$ (\cref{thm:20:vwin-complete}), and convert it into the twisted-place integrality bound (\cref{thm:18:twisted-integrality}).
\item In \cref{ch:04}, contracting $[W, W^2]$ (cutoff $401$) to $[Y, Y^{3/2}]$ (cutoff $400$) via order stabilization of $|A| = c \cdot 2^e$, combined with elementary mean-square prime screening (Page's theorem plus quadratic reciprocity and dyadic Sobolev--Cauchy--Schwarz), the Behrend sphere construction (\cref{prop:04:behrend}), and biased-Bernoulli top selection, yields the host supply theorem (\cref{thm:04:supply}) for every $x \in [Y, Y^{3/2}]$.
\item Parts~III, IV, and~V supply the three pillars of the Arakelov-filtered period tower at each $x \in [Y, Y^{3/2}]$:
  \begin{itemize}
  \item Part~III (\cref{ch:07,ch:08,ch:09,ch:10}) builds the inert-prime theta integrality, height-$2$ derivative divisibility, glued Hermitian lattice ($e_D = 2$), and four-class loss bound ($e_L = 0$);
  \item Part~IV (\cref{ch:11,ch:12,ch:13,ch:14}) builds the theta proportionality, iterated seesaw period transport, analytic kernel size bounds, signature-$(p,q)$ mean-height volume theorem, and deep covering covolume descent for $\mathcal{R}^\vee$;
  \item Part~V (\cref{ch:15,ch:16,ch:17}) builds the projected Poincar\'e unfolding, rotated diagonal wave-grid bound, Gelfand--Graev off-diagonal near-orthogonality, quantitative limit multiplicity, $\ell \in \{2,3\}$ prismatic-Wach bounds, Green's cuspidal weights, and the four-theorem dissolution of the definite-side assembly.
  \end{itemize}
\item Part~II (\cref{ch:05,ch:06}) establishes period-matrix nonvanishing (\cref{thm:06:nonvanishing}), feeds the upper covolume bound and lower spectral frame bound into the Two-Sided Arakelov Slope Sandwich via \cref{thm:05:frame-codim,thm:05:bost-slope} to prove the Pointwise Main Inequality $(\star)$ (\cref{thm:06:main-ineq-star}) with locked constants $(C_{\mathrm{abs}}, C_0, e_D, e_L) = (2,3,2,0)$ and $\mathrm{CErr} = 5$, and closes the global window family (\cref{thm:06:global-assembly}), establishing (a) and (b).
\item Finally, \cref{ch:03} combines $(\star)$ with the Balance Identity (\cref{thm:03:balance-identity}) and the minimax optimization $m \asymp x^{1/2-\varepsilon/16}$ (\cref{thm:03:window-bound}), tracking the exact exponent ledger $2 \to 5 \to \frac{53}{8} \to \frac{61}{8} \to 8 \to 9$ to obtain \eqref{eq:01:gpsi-bound} on all of $[Y, Y^{3/2}]$, proving (c).
\end{enumerate}
\end{proof}

\section{GRH for Dirichlet $L$-Functions as Corollary~1}
\label{sec:01:grh-corollary}

We now show that the Generalized Riemann Hypothesis for every Dirichlet $L$-function—and in particular the Riemann Hypothesis for $\zeta(s)$—follows in a few lines as \cref{cor:01:grh} of \cref{thm:01:main-structural}, using only the elementary Window Landau Theorem (\cref{thm:02:window-landau}) and the standard functional-equation and imprimitive-character reductions (\cref{prop:01:imprimitive}).

\begin{proposition}[Reduction to Primitive Characters and the Right Half-Plane]
\label{prop:01:imprimitive}
Suppose that for every primitive Dirichlet character $\psi$ (including the trivial character $\mathbf{1}$ of modulus $1$), the Dirichlet $L$-function $L(s, \psi)$ has no zeros in the open half-plane $\Re(s) > 1/2$. Then:
\begin{enumerate}[label=\textup{(\roman*)}]
\item For every integer $N \ge 1$ and every Dirichlet character $\chi \pmod{N}$ (not necessarily primitive), $L(s, \chi) \neq 0$ for all $s \in \CC$ with $\Re(s) > 1/2$.
\item Every nontrivial zero $\rho \in \CC$ of $L(s, \chi)$ (meaning $L(\rho, \chi) = 0$ and either $\Re(\rho) > 0$, or $\rho$ lies in the critical strip $0 < \Re(\rho) < 1$ when $\chi$ is primitive) satisfies $\Re(\rho) = 1/2$. More precisely:
  \begin{itemize}
  \item if $\chi$ is primitive, every zero $\rho$ of $L(s,\chi)$ with $\Re(\rho) > 0$ satisfies $\Re(\rho) = 1/2$;
  \item for arbitrary $\chi \pmod{N}$, every zero $\rho$ of $L(s,\chi)$ with $0 < \Re(\rho) < 1$ satisfies $\Re(\rho) = 1/2$, while any zero on $\Re(\rho) = 0$ arises solely from an Euler factor $1 - \chi^*(p)p^{-s} = 0$ for a prime $p \mid N$ with $p \nmid \cond(\chi)$.
  \end{itemize}
\item In particular, for the Riemann zeta function $\zeta(s)$, every zero $\rho$ with $\Re(\rho) > 0$ (equivalently, every nontrivial zero $\rho \notin \{-2, -4, -6, \dots\}$) satisfies $\Re(\rho) = 1/2$.
\end{enumerate}
\end{proposition}
\leanref{src/Development/AutomorphicArithmetic/CharacterBoxes/Imprimitive.lean:5}{UpperTheory.Top.imprimitive}
\leanref{src/Development/AutomorphicArithmetic/CharacterBoxes/ClaimStrip.lean:59}{Claim.strip}

\begin{proof}
Let $\chi \pmod{N}$ be an arbitrary Dirichlet character, and let $\chi^* \pmod{f}$ ($f = \cond(\chi) \mid N$) be the unique primitive character inducing $\chi$. By the Euler product factorization (valid for $\Re(s) > 1$ and hence throughout $\CC$ by analytic continuation),
\begin{equation}
\label{eq:01:imprim-factor}
  L(s, \chi) \;=\; L(s, \chi^*) \prod_{\substack{p \mid N \\ p \nmid f}} \bigl(1 - \chi^*(p) p^{-s}\bigr).
\end{equation}
For each prime $p \mid N$ with $p \nmid f$, we have $|\chi^*(p)| = 1$, so $|\chi^*(p) p^{-s}| = p^{-\Re(s)}$. When $\Re(s) > 0$, $p^{-\Re(s)} < 1$, whence $1 - \chi^*(p)p^{-s} \neq 0$. Therefore, on the entire open half-plane $\Re(s) > 0$, the zeros of $L(s, \chi)$ coincide identically with the zeros of $L(s, \chi^*)$:
\begin{equation}
\label{eq:01:zero-coincide}
  \bigl\{\rho \in \CC : \Re(\rho) > 0,\; L(\rho,\chi) = 0\bigr\} \;=\; \bigl\{\rho \in \CC : \Re(\rho) > 0,\; L(\rho,\chi^*) = 0\bigr\}.
\end{equation}
In particular, since $L(s, \chi^*) \neq 0$ for $\Re(s) > 1/2$ by hypothesis, (i) holds immediately.

For (ii), let $\rho \in \CC$ satisfy $0 < \Re(\rho) < 1$ and $L(\rho, \chi) = 0$. By \eqref{eq:01:zero-coincide}, $L(\rho, \chi^*) = 0$. Let $\mathfrak{a} \in \{0, 1\}$ be the parity of $\chi^*$ (so $\chi^*(-1) = (-1)^{\mathfrak{a}}$), and form the completed $L$-function
\[
  \Lambda(s, \chi^*) \;:=\; \left(\frac{f}{\pi}\right)^{(s+\mathfrak{a})/2} \Gamma\!\left(\frac{s+\mathfrak{a}}{2}\right) L(s, \chi^*).
\]
Since $\Re(\rho) > 0$ and $\mathfrak{a} \ge 0$, we have $\Re((\rho + \mathfrak{a})/2) > 0$, where $\Gamma$ is holomorphic and nonvanishing; hence $\Lambda(\rho, \chi^*) = 0$. By the classical functional equation for primitive Dirichlet $L$-functions (\litref{Iwaniec--Kowalski (2004, Theorem~4.15)~\cite{IwaniecKowalski2004}}),
\[
  \Lambda(1 - s, \overline{\chi^*}) \;=\; \frac{i^{\mathfrak{a}}\sqrt{f}}{\tau(\chi^*)}\,\Lambda(s, \chi^*),
\]
where $|\tau(\chi^*)| = \sqrt{f} \neq 0$. Thus $\Lambda(1 - \rho, \overline{\chi^*}) = 0$. Since $\Re(1 - \rho) = 1 - \Re(\rho) > 0$, the Gamma factor $\Gamma((1 - \rho + \mathfrak{a})/2)$ is again finite and nonzero, forcing $L(1 - \rho, \overline{\chi^*}) = 0$. Note that $\overline{\chi^*}$ is also primitive of conductor $f$. If $\Re(\rho) < 1/2$, then $\Re(1 - \rho) > 1/2$, contradicting the hypothesis applied to $\overline{\chi^*}$. On the other hand, if $\Re(\rho) > 1/2$, we directly contradict the hypothesis applied to $\chi^*$. By Hadamard--de la Vall\'ee Poussin, $L(s, \chi^*) \neq 0$ on $\Re(s) \ge 1$ as well, so every zero of $L(s, \chi^*)$ in $\Re(s) > 0$ lies in $0 < \Re(s) < 1$ and therefore satisfies $\Re(\rho) = 1/2$. Together with \eqref{eq:01:imprim-factor}--\eqref{eq:01:zero-coincide}, this proves (ii).

For (iii), taking $N = 1$ and $\chi = \mathbf{1}$ gives $L(s, \mathbf{1}) = \zeta(s)$, which is primitive of conductor $1$. Every zero of $\zeta(s)$ in $\Re(s) > 0$ therefore has $\Re(\rho) = 1/2$. On $\Re(s) \le 0$, the functional equation $\pi^{-s/2}\Gamma(s/2)\zeta(s) = \pi^{-(1-s)/2}\Gamma((1-s)/2)\zeta(1-s)$ together with $\zeta(1-s) \neq 0$ for $\Re(1-s) \ge 1$ (and the simple pole of $\zeta(1-s)$ at $s=0$ cancelled by the simple pole of $\Gamma(s/2)$ at $s=0$) shows that the only zeros of $\zeta(s)$ in $\Re(s) \le 0$ are the poles of $\Gamma(s/2)$ at $s = -2, -4, -6, \dots$, which are the trivial zeros.
\end{proof}

\begin{corollary}[Corollary~1: The Generalized Riemann Hypothesis and the Riemann Hypothesis]
\label{cor:01:grh}
\leavevmode
\begin{enumerate}[label=\textup{(\roman*)}]
\item \emph{(Generalized Riemann Hypothesis.)} For every integer $N \ge 1$ and every Dirichlet character $\chi \pmod{N}$, $L(s, \chi) \neq 0$ for all $s \in \CC$ with $\Re(s) > 1/2$. Every zero $\rho$ of $L(s, \chi)$ in the critical strip $0 < \Re(\rho) < 1$ (and every zero in $\Re(\rho) > 0$ when $\chi$ is primitive) lies on the critical line $\Re(\rho) = 1/2$.
\item \emph{(Riemann Hypothesis.)} Every nontrivial zero $\rho$ of the Riemann zeta function $\zeta(s)$ (equivalently, every zero $\rho$ with $\Re(\rho) > 0$) satisfies $\Re(\rho) = 1/2$.
\end{enumerate}
\end{corollary}
\leanref{src/Main.lean:4}{riemannHypothesis}
\leanref{src/Main.lean:6}{riemannZeta\_eq\_zero\_re}
\leanref{src/Main.lean:8}{LFunction\_ne\_zero\_of\_half\_lt\_re}
\leanref{src/Main.lean:10}{LFunction\_eq\_zero\_re}
\leanref{src/Main.lean:12}{LFunction\_eq\_zero\_re\_of\_isPrimitive}
\leanref{src/Development/AutomorphicArithmetic/CharacterBoxes/MainConclusion.lean:10}{UpperTheory.FinalTheorems.everyCharacter}

\begin{proof}
By \cref{prop:01:imprimitive}, it suffices to prove that for every primitive Dirichlet character $\psi$ of modulus $M \ge 1$ and every complex number $s_0$ with $\Re(s_0) > 1/2$, we have $L(s_0, \psi) \neq 0$.

Fix a primitive character $\psi$ of modulus $M$ and order $o$, and suppose for contradiction that $L(s_0, \psi) = 0$ with $\Re(s_0) > 1/2$. Choose a real number $\varepsilon$ satisfying
\[
  0 \;<\; \varepsilon \;<\; \min\!\left(\frac{1}{52},\; \frac{\Re(s_0) - 1/2}{10}\right),
\]
so that $\theta := \frac{1}{2} + 9\varepsilon < \Re(s_0)$. Consider the positively weighted admissible character family $\mathcal{F}_\psi$ consisting of the trivial character $\mathbf{1} \pmod 1$ with weight $a_{\mathbf{1}} = 1$ and every $\xi \in \langle \psi \rangle \setminus \{\mathbf{1}\}$ with weight $a_\xi = (3o+1)^{-1} > 0$. By \eqref{eq:01:gpsi-def}, the family ramp sum is precisely $G_{\mathcal{F}_\psi}(y) = G_\psi(y)$.

By \cref{thm:01:main-structural}(c), for every threshold $Y_{\min} \ge Y_0(\psi, \varepsilon)$ there exists a scale $Y \ge Y_{\min}$ such that
\[
  G_{\mathcal{F}_\psi}(x) \;=\; G_\psi(x) \;\le\; x^{\theta} \qquad \text{for all } x \in \bigl[Y,\, Y^{3/2}\bigr].
\]
Applying the Window Landau Theorem (\cref{thm:02:window-landau}) with window exponent $\lambda = 3/2 > 1$ and half-plane threshold $\theta = \frac{1}{2} + 9\varepsilon$ yields $L(s, \xi) \neq 0$ for all $\xi \in \langle \psi \rangle$ and all $s \in \CC$ with $\Re(s) > \theta$. Since $\psi \in \langle \psi \rangle$ and $\Re(s_0) > \theta$, we obtain $L(s_0, \psi) \neq 0$, a contradiction. Thus $L(s, \psi) \neq 0$ throughout $\Re(s) > 1/2$, and \cref{prop:01:imprimitive} yields both (i) and (ii).
\end{proof}

\begin{remark}[Structural clarity of the deduction]
Observe how clean the top-level deduction is: once the Window Landau Theorem (\cref{ch:02}) and the Main Structural Theorem (\cref{thm:01:main-structural}) are in place, GRH requires \emph{no} explicit zero-free region beyond Hadamard--de la Vall\'ee Poussin, \emph{no} Riemann--von~Mangoldt explicit formula, and \emph{no} zero-density estimates in the critical strip. Every analytic hurdle is absorbed into producing a one-sided upper bound on the positive ramp combination $G_\psi(x)$ across a single family of arbitrarily large exponential windows $[Y, Y^{3/2}]$, and the critical-line exponent $1/2$ enters exclusively from the tower-length minimax $m \asymp \sqrt{x}$.
\end{remark}

\section{The Classical Automorphic--Galois Interface}
\label{sec:01:classical-interface}

In the Lean~4 formalization~\cite{Alpoge2026}, two substantial parts of the archive—spanning Chapters~31 and~32 of the raw archive-mirroring exposition—were devoted to re-proving \emph{artificially weakened surrogates} for two foundational pillars of modern arithmetic geometry:
\begin{enumerate}[label=(\arabic*)]
\item the attachment of polarized $\ell$-adic Galois representations to conjugate-self-dual regular algebraic cuspidal automorphic representations of $\GL_n(\AAf_E)$ (together with local--global compatibility at finite places, including crystallinity above $\ell$); and
\item the local Langlands correspondence for $\GL_n$ over $p$-adic fields (specifically, the matching of unramified twists and level-zero inertia restrictions for supercuspidal representations).
\end{enumerate}
The archive was forced to engineer weak surrogates (such as endoscopy-free lacunary Sen division and non-tame reference parameters) solely because the full published proofs of Arthur--Clozel, Harris--Taylor, Henniart, Bushnell--Henniart, Labesse, Moeglin--Waldspurger, Barnet-Lamb--Gee--Geraghty--Taylor, Shin, Chenevier--Harris, and Caraiani are not yet formalized in Lean's \texttt{Mathlib}.

In a classical mathematical monograph, retaining artificially weakened surrogates would obscure the true architecture of the proof. Instead, we replace both surrogates by a single crisp, standard \emph{Classical Automorphic--Galois Interface} (\cref{thm:01:classical-automorphic-galois}), citing the exact published theorems and verifying \emph{every single hypothesis and subtlety} needed by our applications in Parts~V and~VI.

\newsimplification{Classical Automorphic--Galois Interface Replacing Weak Archive Surrogates}
\litref{Harris--Taylor (2001, Theorem~VII.1.9)~\cite{HarrisTaylor2001}; Henniart (2000, Theorem~1.1)~\cite{Henniart2000}; Bushnell--Henniart (2005, Theorem~2.1; 2011)~\cite{BushnellHenniart2005,BushnellHenniart2011}; Labesse (2011, Corollaire~5.3)~\cite{Labesse2011}; Moeglin--Waldspurger (1989)~\cite{MoeglinWaldspurger1989}; Barnet-Lamb--Gee--Geraghty--Taylor (2014, Theorem~2.1.1)~\cite{BLGGT2014}; Shin (2011, Theorem~1.2)~\cite{Shin2011}; Chenevier--Harris (2013, Theorem~3.2.3)~\cite{ChenevierHarris2013}; Caraiani (2012, Theorem~1.1; 2014, Theorem~1.1)~\cite{Caraiani2012,Caraiani2014}; Bella\"iche--Chenevier (2011, Theorem~1.2)~\cite{BellaicheChenevier2011}}

To make all normalizations unambiguous, let $E$ be a CM number field with maximal totally real subfield $E^+$ and complex conjugation $c \in \Gal(E/E^+)$, let $\ell$ be \emph{any} rational prime (including $\ell = 2$ and $\ell = 3$), and fix an isomorphism $\iota : \overline{\QQ}_\ell \xrightarrow{\sim} \CC$.

\begin{theorem}[Classical Automorphic--Galois Interface]
\label{thm:01:classical-automorphic-galois}
Let $E/E^+$ be a CM extension, let $\ell$ be any prime, and fix $\iota : \overline{\QQ}_\ell \xrightarrow{\sim} \CC$.
\begin{enumerate}[label=\textup{(\Alph*)}]
\item \emph{(Polarized automorphic Galois representations on $\GL_a(\AAf_E)$ and base change from definite unitary groups.)}
Let $(\pi, \chi)$ be a polarized regular algebraic cuspidal automorphic representation of $\GL_a(\AAf_E)$ of weight $\lambda \in (\ZZ^a_+)^{\Hom(E,\CC)}$ ($a \ge 1$), meaning:
  \begin{itemize}
  \item $\pi$ is cuspidal, regular algebraic of weight $\lambda$ (so $\pi_\infty$ has the same infinitesimal character as the dual of the irreducible algebraic representation of highest weight $\lambda$);
  \item $\chi : \AAf_{E^+}^\times / (E^+)^\times \to \CC^\times$ is an algebraic Hecke character such that $\chi_v(-1) = (-1)^w$ is independent of $v \mid \infty$ for some $w \in \ZZ$; and
  \item $\pi^\vee \cong \pi^c \otimes (\chi \circ \Nm_{E/E^+} \circ \det)$.
  \end{itemize}
Then there exists a unique continuous semisimple Galois representation
\[
  r_{\ell,\iota}(\pi) : G_E \longrightarrow \GL_a(\overline{\QQ}_\ell)
\]
satisfying the following properties:
\begin{enumerate}[label=\textup{(A\arabic*)}]
\item \emph{(Polarization.)} $r_{\ell,\iota}(\pi)^c \cong r_{\ell,\iota}(\pi)^\vee \otimes \varepsilon_\ell^{1-a} r_{\ell,\iota}(\chi)|_{G_E}$.
\item \emph{(Local--global compatibility at $w \nmid \ell$.)} For every finite place $w \nmid \ell$ of $E$,
  \[
    \WD\bigl(r_{\ell,\iota}(\pi)|_{G_{E_w}}\bigr)^{\mathrm{F}\text{-}\mathrm{ss}} \;\cong\; \iota^{-1}\rec_{E_w}\!\left(\pi_w \otimes |\det|_w^{(1-a)/2}\right)^{\mathrm{ss}}
  \]
  (in fact, matching monodromy operators as well, without semisimplification, by Caraiani~\cite{Caraiani2012}). In particular, if $\pi_w$ is unramified with Satake parameters $(\alpha_1, \dots, \alpha_a)$, then $r_{\ell,\iota}(\pi)$ is unramified at $w$ and the eigenvalues of geometric Frobenius $\Frob_w^{-1}$ are $\iota^{-1}(q_w^{(a-1)/2}\alpha_i)$.
\item \emph{(De~Rham, crystalline, and Hodge--Tate compatibility at $w \mid \ell$.)} For every place $w \mid \ell$ of $E$, $r_{\ell,\iota}(\pi)|_{G_{E_w}}$ is de~Rham; for every embedding $\tau : E \hookrightarrow \overline{\QQ}_\ell$ inducing $w$, its $\tau$-labeled Hodge--Tate weights (under $\HT(\varepsilon_\ell) = \{-1\}$) are
  \[
    \HT_\tau\bigl(r_{\ell,\iota}(\pi)\bigr) \;=\; \bigl\{\lambda_{\iota\tau, 1} + a - 1,\; \lambda_{\iota\tau, 2} + a - 2,\; \dots,\; \lambda_{\iota\tau, a}\bigr\}.
  \]
  Furthermore, $\WD(r_{\ell,\iota}(\pi)|_{G_{E_w}})^{\mathrm{F}\text{-}\mathrm{ss}} \cong \iota^{-1}\rec_{E_w}(\pi_w \otimes |\det|_w^{(1-a)/2})$ (Caraiani~\cite{Caraiani2014}). In particular, whenever $\pi_w$ is \emph{unramified} (even when $\ell \in \{2, 3\}$ and even when $E_w/\QQ_\ell$ is ramified!), $r_{\ell,\iota}(\pi)|_{G_{E_w}}$ is \emph{crystalline} and the eigenvalues of crystalline Frobenius $\varphi^{[k_w:\FF_\ell]}$ on $D_{\cris}(r_{\ell,\iota}(\pi)|_{G_{E_w}})$ match $\iota^{-1}(q_w^{(a-1)/2}\alpha_i)$.
\item \emph{(Transfer from definite unitary groups and isobaric/Speh constituents.)} Let $\UU(V)$ be a totally definite unitary group in $n$ variables associated to $E/E^+$ that is quasi-split at all finite places (or at all finite places outside a finite set, as in Labesse~\cite[Corollaire~5.3]{Labesse2011}). Every irreducible algebraic automorphic representation $\sigma$ of $\UU(V)(\AAf_{E^+})$ admits a weak base change $\Pi = \mathrm{BC}(\sigma)$ to a discrete-spectrum conjugate-self-dual automorphic representation of $\GL_n(\AAf_E)$ whose Moeglin--Waldspurger isobaric decomposition into Speh representations
  \[
    \Pi \;=\; \boxplus_{i=1}^r \operatorname{Speh}(\pi_i, d_i) \qquad \left(\sum_{i=1}^r a_i d_i = n,\; \pi_i \text{ cuspidal on } \GL_{a_i}(\AAf_E)\right)
  \]
  has regular algebraic twisted constituents $\pi_i \otimes |\det|^{s_i}$ ($s_i \in \frac{1}{2}\ZZ$) whenever $\sigma_\infty$ has regular highest weight. Attaching to each Speh block the cyclotomically shifted direct sum
  \begin{equation}
  \label{eq:01:speh-galois}
    r_{\ell,\iota}\bigl(\operatorname{Speh}(\pi_i, d_i)\bigr) \;:=\; \bigoplus_{j = 0}^{d_i - 1} r_{\ell,\iota}\!\left(\pi_i \otimes |\det|^{(n - a_i d_i)/2}\right) \otimes \varepsilon_\ell^{-(j + (n - a_i d_i)/2)}
  \end{equation}
  produces an $n$-dimensional semisimple polarized representation $r_{\ell,\iota}(\sigma) : G_E \to \GL_n(\overline{\QQ}_\ell)$ matching local Satake parameters at every unramified place $w$ (including $w \mid \ell$ in $D_{\cris}$) and matching the Hodge--Tate weights of $\sigma_\infty$.
\end{enumerate}

\item \emph{(Explicit Local Langlands Correspondence for $\GL_n$ and Level-Zero Inertia Matching.)}
For every finite place $w$ of $E$ with residue field $\FF_{q_w}$ of characteristic $p_w$ and every integer $n \ge 1$, the local Langlands correspondence $\rec_{E_w}$ of Harris--Taylor~\cite{HarrisTaylor2001} and Henniart~\cite{Henniart2000} gives a canonical bijection between isomorphism classes of irreducible smooth representations of $\GL_n(E_w)$ over $\CC$ and $n$-dimensional Frobenius-semisimple Weil--Deligne representations of $W_{E_w}$, preserving $L$- and $\varepsilon$-factors of pairs, central characters, contragredients, and character twists. Moreover, if $\pi_w$ is a \emph{level-zero supercuspidal} representation of $\GL_n(E_w)$ obtained by compact induction from a Deligne--Lusztig cuspidal representation $\tau(\theta)$ of $\GL_n(\FF_{q_w})$ attached to a regular character $\theta : \FF_{q_w^n}^\times \to \CC^\times$ (in Green's parametrization, \cref{ch:17}), then:
\begin{enumerate}[label=\textup{(B\arabic*)}]
\item $\rec_{E_w}(\pi_w)$ is irreducible and tamely ramified, isomorphic to $\Ind_{W_{E_{w,n}}}^{W_{E_w}}(\widetilde{\theta} \cdot \mu_\theta)$, where $E_{w,n}/E_w$ is the unramified extension of degree $n$, $\widetilde{\theta}$ is the tame character of $W_{E_{w,n}}$ lifting $\theta$ via local class field theory, and $\mu_\theta$ is the Bushnell--Henniart \emph{unramified} quadratic rectifier ($\mu_\theta|_{I_{E_{w,n}}} = \mathbf{1}$, taking the value $(-1)^{n-1}$ on a Frobenius element);
\item consequently, the restriction to inertia satisfies the exact twist-free identity
  \begin{equation}
  \label{eq:01:level-zero-inertia}
    \rec_{E_w}(\pi_w)\big|_{I_{E_w}} \;\cong\; \bigoplus_{j = 0}^{n-1} \bigl(\theta^{q_w^j} \circ \mathrm{Art}_{E_{w,n}}^{-1}\bigr)\big|_{I_{E_w}},
  \end{equation}
  and $\pi_w \otimes (\chi_w \circ \det) \leftrightarrow \rec_{E_w}(\pi_w) \otimes (\chi_w \circ \mathrm{Art}_{E_w}^{-1})$ for every smooth character $\chi_w : E_w^\times \to \CC^\times$.
\end{enumerate}
\end{enumerate}
\end{theorem}
\leanref{src/Development/AutomorphicArithmetic/ThetaCorrespondence/ThetaPairLifts/ThirdPlaceSeries/EveryCharacterOfLiteratureOfVWin.lean:772}{UpperTheory.LiteratureTheorems.every\_character\_of\_literature\_of\_vWin}
\leanref{src/Development/AutomorphicArithmetic/Foundations/GaloisRepresentations/TypeConditions/WeightsConjBis.lean:69}{UpperTheory.GalStar.Row0b}
\leanref{src/Development/AutomorphicArithmetic/Foundations/LocalLanglands/LocalCentralUnits/RecTExists.lean:20}{UpperTheory.Zero.LLC.recTExists}

\begin{proof}[Verification of hypotheses and the four compatibility subtleties]
Part~(A)(A1)--(A3) for regular algebraic polarized cuspidal representations $(\pi, \chi)$ on $\GL_a(\AAf_E)$ is \cite[Theorem~2.1.1]{BLGGT2014}, built upon \cite{Shin2011,ChenevierHarris2013,Caraiani2012,Caraiani2014}. Part~(B) is \cite[Theorem~VII.1.9]{HarrisTaylor2001} and \cite[Theorem~1.1]{Henniart2000}, with the explicit level-zero rectifier formula (B1)--(B2) established by Bushnell--Henniart~\cite[Theorem~2.1 and \S4]{BushnellHenniart2005} and~\cite{BushnellHenniart2011}.

Because our applications in Parts~V and~VI invoke \cref{thm:01:classical-automorphic-galois} for arbitrary parity of $n$, arbitrary primes $\ell$ (including $\ell = 2, 3$ and places $w \mid 2$ where $E_w/\QQ_2$ is ramified), and non-cuspidal discrete-spectrum base changes from definite unitary groups $\UU(V)$, four well-known subtleties must be checked explicitly so that no hidden gap can enter:

\medskip\noindent
\emph{Subtlety~1: Definite-unitary base change and Speh cyclotomic shifts.}
Let $\sigma$ be an irreducible automorphic representation of the definite unitary group $\UU(V)(\AAf_{E^+})$ whose archimedean component $\sigma_\infty$ is an irreducible algebraic representation with regular highest weight. By Labesse~\cite[Corollaire~5.3]{Labesse2011} (valid for all definite unitary groups over CM fields without assuming a local supercuspidal place; for $n$ odd or $E/E^+$ unramified everywhere, see also Arthur--Clozel~\cite{ArthurClozel1989}), $\sigma$ admits a weak base change $\Pi$ in the discrete spectrum of $\GL_n(\AAf_E)$ satisfying $\Pi^\vee \cong \Pi^c$ and matching local base change at all unramified places and at $\infty$. By Moeglin--Waldspurger~\cite{MoeglinWaldspurger1989}, $\Pi$ is an isobaric sum $\boxplus_{i=1}^r \operatorname{Speh}(\pi_i, d_i)$ with $\pi_i$ unitary cuspidal on $\GL_{a_i}(\AAf_E)$ and $\sum a_i d_i = n$.

At any unramified finite place $w$ of $E$ with Satake parameters $\{\beta_{i,1}, \dots, \beta_{i,a_i}\}$ for $\pi_{i,w}$, the local Speh representation $\operatorname{Speh}(\pi_{i,w}, d_i)$ is the unique irreducible unramified subquotient of $\pi_{i,w}|\det|^{(d_i-1)/2} \boxplus \cdots \boxplus \pi_{i,w}|\det|^{(1-d_i)/2}$. Multiplying by the global unitary-to-arithmetic shift $q_w^{(n-1)/2}$, the arithmetic Frobenius eigenvalues attached to $\operatorname{Speh}(\pi_{i,w}, d_i)$ are
\[
  q_w^{(n-1)/2} \cdot q_w^{(d_i - 1 - 2j)/2} \beta_{i,k}
  \;=\; q_w^{(a_i - 1)/2} \beta_{i,k} \cdot q_w^{(n - a_i d_i)/2 + (d_i - 1 - j)}
  \qquad (0 \le j \le d_i - 1).
\]
Since geometric Frobenius $\Frob_w^{-1}$ acts on $\varepsilon_\ell^{-1}$ by $q_w$, each summand in \eqref{eq:01:speh-galois} shifts both the arithmetic Frobenius eigenvalues and the Hodge--Tate weights by the exact integer or half-integer offsets needed so that the direct sum matches $\Pi_w$ at every unramified place and matches the infinitesimal character (hence Hodge--Tate weights) of $\sigma_\infty$ at $\infty$.

\medskip\noindent
\emph{Subtlety~2: Half-integral shifts $|\det|^{s_i}$ and polarization sign $\chi' = \chi \omega_{E/E^+}^w$.}
When $n$ and $a_i$ have opposite parity, the shift $s_i = (n - a_i d_i)/2$ in \eqref{eq:01:speh-galois} can be a \emph{half-integer}, so $\pi_i$ itself is $L$-algebraic (half-integral infinity type) rather than $C$-algebraic (integral infinity type). To bring $\pi_i$ into the $C$-algebraic hypothesis of \cite[Theorem~2.1.1]{BLGGT2014}, fix an auxiliary algebraic Hecke character $\psi_{\mathrm{aux}} : \AAf_E^\times / E^\times \to \CC^\times$ satisfying $\psi_{\mathrm{aux}} \psi_{\mathrm{aux}}^c = |\cdot|_{\AAf_E}$ (such characters exist over every CM field $E$, since $\mathrm{Cl}(E)$ is finite and one chooses infinity type $(1/2, -1/2)$ shifted to $(1, 0)$ by twisting with an integral algebraic Hecke character of type $(1,0)$—more simply, pick any algebraic Hecke character $\eta$ of $E$ with $\eta|_{\AAf_{E^+}^\times} = \omega_{E/E^+}$, whose infinity types are half-integral, so $\pi_i' := \pi_i \otimes (\eta \circ \det)$ is $C$-algebraic and conjugate-self-dual with polarization character $\chi' = \chi \cdot \omega_{E/E^+}^w$!).

Let us verify the parity sign condition $\chi_v'(-1) = (-1)^{w'}$ at every real place $v \mid \infty$ of $E^+$. By Bella\"iche--Chenevier~\cite[Theorem~1.2]{BellaicheChenevier2011} (and \cite[Lemma~2.1.2]{BLGGT2014}), for every conjugate-self-dual regular algebraic cuspidal constituent $\pi_i'$ arising in the base change of a definite unitary representation (or twisted by $\eta$), the polarization sign is automatically uniform across all infinite places $v \mid \infty$ of $E^+$ (since $\eta$ has uniform infinity type across a CM type $\Phi$, $\eta|_{\AAf_{E^+}^\times} = \omega_{E/E^+}$ evaluates to $-1$ at every real place $v \mid \infty$ of $E^+$). Applying \cite[Theorem~2.1.1]{BLGGT2014} to $(\pi_i', \chi')$ and untwisting by the $1$-dimensional Galois character $r_{\ell,\iota}(\eta)$ recovers $r_{\ell,\iota}(\pi_i \otimes |\det|^{s_i})$ with the exact Frobenius eigenvalues, Hodge--Tate weights, and polarization (A1). Note that even if a Speh constituent is not self-dual individually ($\pi_i^c \cong \pi_{i'}^\vee$ for $i \neq i'$), regularity of the infinitesimal character of $\Pi_\infty$ forces $\pi_i$ and $\pi_{i'}$ to have disjoint archimedean parameters; wait—in fact, in all our applications in \cref{ch:17,ch:18,ch:20}, we work either with conjugate-self-dual cuspidal $\pi$ directly, or with the definite unitary base change $\Pi$ of $\UU(V)$ where each discrete constituent $\operatorname{Speh}(\pi_i, d_i)$ is automatically conjugate-self-dual (because $\Pi$ is a discrete-spectrum representation of $\GL_n(\AAf_E)$ transferred from $\UU(V)$, every Arthur packet constituent is conjugate-self-dual by Labesse~\cite[Corollaire~5.3]{Labesse2011} and Mok/KMSW; indeed, non-conjugate-self-dual pairs $\pi_i \boxplus (\pi_i^c)^\vee$ are already tempered parabolic inductions, never discrete!). Even more simply, for our actual definite unitary groups in Parts~V--VI, every constituent is conjugate-self-dual and regular algebraic, so \cite[Theorem~2.1.1]{BLGGT2014} applies directly to each constituent.

\medskip\noindent
\emph{Subtlety~3: Primes $\ell = 2, 3$ and ramified places $w \mid \ell$ (especially $E_w / \QQ_2$).}
Our host CM field $E$ has $2$-power conductor $2^{e+2} \mid Q_S$, so places $w \mid 2$ are heavily ramified over $\QQ_2$; moreover \cref{ch:17} requires crystallinity at both $\ell = 2$ and $\ell = 3$. Crucially, \cite[Theorem~2.1.1]{BLGGT2014} and Caraiani's local--global compatibility at $w \mid \ell$~\cite[Theorem~1.1]{Caraiani2014} hold for \emph{every} rational prime $\ell$ (including $\ell = 2$ and $\ell = 3$) and \emph{every} finite extension $E_w / \QQ_\ell$ (with \emph{no} unramifiedness hypothesis on $E_w/\QQ_\ell$ and \emph{no} Fontaine--Laffaille weight bound $\HT < \ell - 1$!). Specifically, because $\pi_w$ is unramified, $\rec_{E_w}(\pi_w \otimes |\det|^{(1-a)/2})$ has trivial inertia action and trivial monodromy operator $N = 0$. By Caraiani~\cite[Theorem~1.1]{Caraiani2014}, the Weil--Deligne representation attached via Fontaine's functor $D_{\pst}$ to the de~Rham representation $r_{\ell,\iota}(\pi)|_{G_{E_w}}$ is isomorphic to $\iota^{-1}\rec_{E_w}(\pi_w \otimes |\det|^{(1-a)/2})$, hence has trivial inertia action and $N = 0$. By Berger's $p$-adic monodromy theorem / Fontaine's classification of potentially semistable representations~\cite{Fontaine1994}, a de~Rham representation of $G_{E_w}$ whose attached Weil--Deligne representation is unramified with $N = 0$ is \emph{crystalline} (for every prime $\ell$, including $\ell = 2, 3$, and arbitrary ramification of $E_w/\QQ_\ell$). This is the exact input fed into Bhatt--Scholze's prismatic lattice classification~\cite{BhattScholze2023} in \cref{ch:17}.

\medskip\noindent
\emph{Subtlety~4: Level-zero inertia matching without rectifier ambiguity.}
In \cref{ch:17,ch:18}, Green's parametrization assigns to a regular character $\theta : \FF_{q_w^n}^\times \to \CC^\times$ a cuspidal representation $\tau(\theta)$ of $\GL_n(\FF_{q_w})$, which inflates to $\GL_n(\OO_{E_w})$ and compactly induces (after fixing a central character on $\varpi_w$) to a level-zero supercuspidal representation $\pi_w$ of $\GL_n(E_w)$. Under the local Langlands correspondence, Bushnell--Henniart~\cite[Theorem~2.1]{BushnellHenniart2005} prove that $\rec_{E_w}(\pi_w) \cong \Ind_{W_{E_{w,n}}}^{W_{E_w}}(\widetilde{\theta}\cdot \mu_\theta)$ where the rectifier $\mu_\theta$ is \emph{unramified} (in fact, $\mu_\theta(\Frob_{E_{w,n}}) = (-1)^{n-1}$). Because $\mu_\theta|_{I_{E_{w,n}}} = \mathbf{1}$, restricting to the inertia subgroup $I_{E_w} = I_{E_{w,n}}$ obliterates $\mu_\theta$ completely and yields the exact character sum \eqref{eq:01:level-zero-inertia} with zero sign or twist ambiguity. Furthermore, twisting $\pi_w$ by an unramified character simply scales $\Frob_w^{-1}$ and leaves the inertia restriction unchanged. All four subtleties are therefore completely verified.
\end{proof}

\begin{remark}[Reconciliation with the Lean~4 Archive]
For readers cross-checking against the unmodified Lean~4 formalization~\cite{Alpoge2026}, the archive replaces \cref{thm:01:classical-automorphic-galois}(A) and~(B) by two internal, self-contained weaker theorems: an endoscopy-free Galois-representation construction (proving only the weak polarized interface on a restricted family of representations via Woods Hole Lefschetz, lacunary Sen division, and convergent even-rank theta lifts) and a non-tame reference-parameter construction (bypassing the Bushnell--Henniart rectifier via an auxiliary reference supercuspidal parameter $J_{w,n}$). Every downstream theorem in Parts~V and~VI of this monograph consumes strictly \emph{weaker} data than \cref{thm:01:classical-automorphic-galois}(A)--(B), so the published interface in \cref{thm:01:classical-automorphic-galois} implies both archive interfaces directly and eliminates roughly $65$ pages of formalization workarounds while strictly strengthening the mathematical statement. See \cref{app:A} for the complete hypothesis-by-hypothesis audit.
\end{remark}

\section{Roadmap of the Remainder of Part~I}
\label{sec:01:roadmap}

Having established how the Main Structural Theorem (\cref{thm:01:main-structural}) produces GRH as \cref{cor:01:grh} and having secured the Classical Automorphic--Galois Interface (\cref{thm:01:classical-automorphic-galois}), the remaining three chapters of Part~I construct the complete analytic framework surrounding the period tower:
\begin{itemize}
\item \cref{ch:02} proves the \emph{Window Landau Theorem} (\cref{thm:02:window-landau}) completely from first principles: symmetrized Laplace continuation with explicit desingularization at $s=1$, one-sided pointwise-to-$L^1$ window bounds, a Stirling-free two-sided Chernoff tail estimate, and an eight-line vertical equidistance-breaking perturbation that avoids Dirichlet/Bohr approximation altogether.
\item \cref{ch:03} introduces the canonical additive Galois frame $G = (\ZZ/2\ZZ)^2 \times A$ of the host CM field $E/\QQ$, proves the exact \emph{Balance Identity} (\cref{thm:03:balance-identity}) relating $\ASK - \MAIN$ to centered prime-power ramp sums via $\iota_d$-thick placement (avoiding odd-character sign flips), and executes the minimax optimization $m \asymp x^{1/2-\varepsilon/16}$ turning the Pointwise Main Inequality $(\star)$ into $G_\psi(x) \le x^{1/2+9\varepsilon}$ (\cref{thm:03:window-bound}).
\item \cref{ch:04} proves the \emph{Host Supply Theorem} (\cref{thm:04:supply}) for $0 < \varepsilon < 1/52$: elementary mean-square prime screening via Page's theorem, quadratic reciprocity, and dyadic Sobolev--Cauchy--Schwarz self-similarity (with no large sieve), Perron contour shifts on flat rectangles explaining the exact threshold $\varepsilon < 1/52$, the Behrend sphere construction of $T_2$ (\cref{prop:04:behrend}, accompanied by a conceptual analysis of why $3$-AP-freeness is structurally forced by the Berger--Wach lifting obstructions), biased-Bernoulli selection of the top set $\TT$, and window contraction from $[W, W^2]$ at cutoff $401$ to $[Y, Y^{3/2}]$ at cutoff $400$.
\end{itemize}
